\documentclass[11pt]{article}
\usepackage[letterpaper,margin=0.9in]{geometry}
\usepackage{amsmath,amssymb,amsthm,booktabs,microtype,xurl}
\usepackage[hidelinks]{hyperref}
\hypersetup{
  pdftitle={Eternal domination and the Lovasz theta function: an unbounded separation and a smallest counterexample},
  pdfauthor={Alec Kriebel},
  pdfsubject={An unbounded separation and an exact smallest counterexample},
  pdfkeywords={eternal domination, Lovasz theta function, graph theory, semidefinite programming}
}

\newtheorem{theorem}{Theorem}
\newtheorem{corollary}[theorem]{Corollary}
\newcommand{\cc}{\operatorname{cc}}

\title{Eternal domination and the Lov\'asz theta function:\\
an unbounded separation and a smallest counterexample}
\author{Alec Kriebel\\[2pt]\large with heavy assistance from ChatGPT 5.6 Sol}
\date{First prepared 25 July 2026\\Revised 26 July 2026\\Not peer reviewed}

\begin{document}
\maketitle

\begin{abstract}
We show that the Lov\'asz theta function can exceed the standard one-guard
eternal domination number by an unbounded factor.  Combining results of Alon
and of Goddard, Hedetniemi, and Hedetniemi gives an explicit infinite family
\((H_k)\) with
\[
 2\leq\gamma^\infty(H_k)\leq3
 \quad\text{and}\quad
 \vartheta(H_k)=\Theta\!\left(|V(H_k)|^{1/3}\right).
\]
Thus the ratio \(\vartheta/\gamma^\infty\) is unbounded.  Second, let \(G\)
be the ten-vertex graph with graph6 record \texttt{IEhbtj\{ro}.
MacGillivray, Mynhardt, and Virgile proved that
\(\gamma^\infty(G)=3\).  We give an exact rational feasible matrix for the
standard Lov\'asz-theta semidefinite program with objective
\[
 \frac{7593}{2500}=3.0372.
\]
Their exhaustive result through order nine, together with the theta sandwich
inequality, shows that \(G\) has minimum possible order among counterexamples.
A solver-free verifier checks the finite theta certificate and independently
recomputes the eternal-domination game.
\end{abstract}

\noindent\textbf{Verification disclaimer.}
This is an unreviewed AI-assisted research note released for independent
checking.  Exact computation verifies the encoded finite statements; it is
not peer review and does not establish literature priority.

\section{Introduction and conventions}

In the one-guard eternal domination game on a finite simple graph \(G\), at
most one guard occupies each vertex and the occupied vertices form a
dominating set.  At each turn an unoccupied vertex is attacked, and exactly
one guard must move along an edge to the attacked vertex so that the new guard
set is again dominating.  The minimum number of guards that can defend
indefinitely is \(\gamma^\infty(G)\)~\cite{burger}.
The elementary bounds for this game are
\begin{equation}\label{eq:guard-bounds}
 \alpha(G)\leq\gamma^\infty(G)\leq\cc(G),
\end{equation}
where \(\cc(G)\) denotes the clique-cover number~\cite[Theorem~1]{ghh}.

West's problem page asks whether \(\gamma^\infty\) is bounded below by the
Lov\'asz theta function~\cite{west}.  There is a potential notation trap.
Work on eternal domination often writes \(\theta(G)\) for the clique-cover
number.  Here we instead write
\[
 \cc(G)=\chi(\overline G)
\]
for clique cover and reserve \(\vartheta(G)\) for the Lov\'asz theta function.
Lov\'asz introduced the function in 1979~\cite{lovasz}.  With the standard
convention used by Alon and Kahale~\cite{alon-kahale},
\begin{equation}\label{eq:theta}
 \vartheta(G)=
 \max\left\{\langle J,X\rangle:
 X\succeq0,\ \operatorname{Tr}(X)=1,\
 X_{ij}=0\text{ for }ij\in E(G)\right\}.
\end{equation}
It satisfies the familiar sandwich
\begin{equation}\label{eq:sandwich}
 \alpha(G)\leq\vartheta(G)\leq\chi(\overline G)=\cc(G).
\end{equation}

We give two complementary forms of the negative answer.  Section~2 records
an unbounded polynomial separation that follows directly from established
results.  Section~3 gives an exact smallest counterexample, together with a
solver-free certificate.

\section{An explicit unbounded separation}

Alon's explicit Ramsey construction and a theorem on eternal security combine
to give considerably more than a single counterexample.

\begin{theorem}\label{thm:unbounded}
For every integer \(k\geq2\) not divisible by three, put \(n_k=2^{3k}\).
There is an explicit graph \(H_k\) on \(n_k\) vertices such that
\[
 \alpha(H_k)=2,\qquad
 2\leq\gamma^\infty(H_k)\leq3,
 \qquad
 \vartheta(H_k)=\Theta(n_k^{1/3}),
\]
and hence
\[
 \frac{\vartheta(H_k)}{\gamma^\infty(H_k)}
 =\Theta(n_k^{1/3}).
\]
In particular, \(\vartheta/\gamma^\infty\) is unbounded, and there is no
constant \(C\) such that
\(\vartheta(H)\leq C\gamma^\infty(H)\) for every graph \(H\).
\end{theorem}

\begin{proof}
Alon explicitly constructs a
\(2^{k-1}(2^{k-1}-1)\)-regular triangle-free Cayley graph \(F_k\) on
\(n_k\) vertices whose complement satisfies
\[
 \vartheta(\overline{F_k})=\Theta(n_k^{1/3})
\]
as \(k\) ranges over these values~\cite[Corollary~2.3]{alon-ramsey}.  Set
\(H_k=\overline{F_k}\).  Since \(k\geq2\), the graph \(F_k\) has an edge,
and triangle-freeness gives
\[
 \alpha(H_k)=\omega(F_k)=2.
\]
Goddard, Hedetniemi, and Hedetniemi proved that every graph with independence
number two has one-guard eternal domination number at most
three~\cite[Theorem~4]{ghh}.  Their notation \(\sigma_1\) denotes the
parameter written \(\gamma^\infty\) here.  Equation~\eqref{eq:guard-bounds}
therefore gives
\[
 2\leq\gamma^\infty(H_k)\leq3.
\]
Dividing Alon's theta estimate by a quantity bounded between two and three
gives the asserted ratio.
\end{proof}

\section{A smallest counterexample}

MacGillivray, Mynhardt, and Virgile found two smallest graphs for which
\(\gamma^\infty<\cc\)~\cite{mmv}.  Their Table~9 records the graph6 strings
\texttt{IEhbtj\{ro} and \texttt{IEhbtn\{ro}.  Their Fact~5.1 gives
clique-cover number four, while Proposition~5.1 and its proof give eternal
domination number three.  The observation here is that the 26-edge member of
that pair is already a counterexample to the proposed Lov\'asz-theta lower
bound.

Let \(G\) be the graph on vertices \(0,\ldots,9\) whose graph6 record is
\[
 \texttt{IEhbtj\{ro}.
\]
For complete independence from graph6 software, its edge set is
\begin{align*}
E(G)=\{&
03,04,07,08,09,\,
13,15,16,18,19,\,
24,25,26,27,28,\\
&36,37,38,\,
46,48,49,\,
57,58,59,\,
69,79\},
\end{align*}
where \(ij\) denotes the unordered edge \(\{i,j\}\).

\begin{theorem}\label{thm:counterexample}
The graph \(G\) satisfies
\[
 \gamma^\infty(G)=3
 <\frac{7593}{2500}
 \leq\vartheta(G).
\]
Consequently, the inequality
\(\gamma^\infty(H)\geq\vartheta(H)\) does not hold for all graphs \(H\).
\end{theorem}

\begin{proof}
The equality \(\gamma^\infty(G)=3\) is established in Proposition~5.1 and
its proof by MacGillivray, Mynhardt, and Virgile~\cite{mmv}.  We also
independently verify it by exact state enumeration in
Section~\ref{sec:verification}.

For the theta bound, define \(X=M/10000\), where
\begingroup
\scriptsize
\[
M=\begin{pmatrix}
1170&576&621&0&0&576&621&0&0&0\\
576&1172&621&0&576&0&0&621&0&0\\
621&621&1025&335&0&0&0&0&0&512\\
0&0&335&1025&621&621&0&0&0&512\\
0&576&0&621&1172&576&0&621&0&0\\
576&0&0&621&576&1172&621&0&0&0\\
621&0&0&0&0&621&1025&335&512&0\\
0&621&0&0&621&0&335&1025&512&0\\
0&0&0&0&0&0&512&512&607&196\\
0&0&512&512&0&0&0&0&196&607
\end{pmatrix}.
\]
\endgroup
The matrix is symmetric, its diagonal entries sum to \(10000\), and direct
inspection against the displayed edge list gives \(M_{ij}=0\) for every
\(ij\in E(G)\).

It remains to verify positive semidefiniteness without numerical
approximation.  Exact Gaussian elimination gives the following diagonal
pivots in an \(LDL^{\mathsf T}\) factorization of \(M\):
\[
\begin{gathered}
1170,\quad
\frac{57748}{65},\quad
\frac{33696545}{57748},\quad
\frac{5611637865}{6739309},\quad
\frac{5875307251}{81328085},\\
\frac{147954747608}{2670594205},\quad
\frac{774768046339541}{1627502223688},\quad
\frac{47387055721842240}{774768046339541},\\
\frac{1491051941429}{39596916715},\quad
\frac{50290089836283}{1491051941429}.
\end{gathered}
\]
Every pivot is positive, so \(M\), and hence \(X\), is positive definite.
Thus \(X\) is feasible in~\eqref{eq:theta}.  Finally, the sum of all entries
of \(M\) is \(30372\), and therefore
\[
 \langle J,X\rangle
 =\frac{30372}{10000}
 =\frac{7593}{2500}
 >3.
\]
Feasibility gives the asserted lower bound for \(\vartheta(G)\), and the
strict separation follows.
\end{proof}

\begin{corollary}\label{cor:minimal}
The minimum order of a graph \(H\) satisfying
\(\gamma^\infty(H)<\vartheta(H)\) is ten.
\end{corollary}

\begin{proof}
MacGillivray, Mynhardt, and Virgile computationally verified that every graph
of order at most nine satisfies
\(\gamma^\infty(H)=\cc(H)\)~\cite[Proposition~5.2]{mmv}.  By
\eqref{eq:sandwich}, every such graph satisfies
\[
 \vartheta(H)\leq\cc(H)=\gamma^\infty(H).
\]
Theorem~\ref{thm:counterexample} supplies a counterexample of order ten.
\end{proof}

\section{Independent exact verification}\label{sec:verification}

The accompanying verifier uses Python's standard library and no SDP solver.
It parses the graph6 record independently in two implementations and compares
the result with the frozen edge list above.  For the theta certificate it
checks symmetry, trace, all 26 edge zeros, the exact entry sum, and the ten
rational \(LDL^{\mathsf T}\) pivots.

For the eternal-domination game, let \(\mathcal D_k\) be the dominating
\(k\)-subsets of \(V(G)\).  Starting with
\(\mathcal F_0=\mathcal D_k\), repeatedly set
\[
\mathcal F_{t+1}=
\left\{S\in\mathcal F_t:
\begin{array}{l}
\text{for every }v\notin S\text{ there is }u\in S\cap N(v)\\
\text{such that }(S\setminus\{u\})\cup\{v\}\in\mathcal F_t
\end{array}
\right\}.
\]
Because the state space is finite, this deletion process reaches the greatest
closed family.  A nonempty closed family is equivalent to an eternal defense:
it supplies a response after every attack, while the reachable states of any
successful defense form such a family.

The exact fixed-point sizes for \(k=1,2,3\) are
\[
 0,\quad 0,\quad 86.
\]
There are 86 dominating triples, all of which survive the fixed point.
The saved certificate lists them and the verifier checks a legal successor
inside the family for all \(86(10-3)=602\) configuration-attack pairs.  This
independently reproduces \(\gamma^\infty(G)=3\).

\section{Attribution, novelty, and scope}

The ingredients of Theorem~\ref{thm:unbounded} are established results:
Alon supplies the explicit theta construction, and Goddard, Hedetniemi, and
Hedetniemi supply the constant eternal-domination bound.  The theorem is a
direct synthesis, not a new graph construction.  Likewise, the ten-vertex
graph and its eternal-domination value are due to MacGillivray, Mynhardt, and
Virgile~\cite{mmv}.  The candidate contribution is recognizing that these
results answer West's Lov\'asz-theta question, together with the exact
rational certificate and minimum-order identification.

A focused search through notation variants, the cited papers' reference and
citation graphs, recent related papers, public problem pages, arXiv, and
scholarly indexes found no prior public source explicitly making either
eternal-domination connection as of 26 July 2026.  That search cannot
establish worldwide priority: an unindexed, unpublished, or differently
phrased antecedent may exist.

The counterexample theorem is an exact finite result.  The minimum-order
corollary additionally relies on the published exhaustive computation through
order nine.  The unbounded-family theorem relies on the cited published
results.  No numerical estimate of the optimum theta value of the ten-vertex
graph is needed.

\section*{AI-assistance and authorship disclosure}

The observation, exact certificates, independent verifiers, literature
review, proof drafting, and publication materials were developed with heavy
assistance from ChatGPT 5.6 Sol under Alec Kriebel's direction.  Alec Kriebel
is the named human author.  No external expert was contacted or has reviewed
this note.  Passing the supplied checks is not peer review.

\clearpage
\begin{thebibliography}{9}

\bibitem{alon-ramsey}
N.~Alon,
\newblock Explicit Ramsey graphs and orthonormal labelings,
\newblock \emph{Electronic Journal of Combinatorics} 1 (1994),
Research Paper 12, 8 pages.
\newblock \url{https://doi.org/10.37236/1192}

\bibitem{alon-kahale}
N.~Alon and N.~Kahale,
\newblock Approximating the independence number via the
\(\theta\)-function,
\newblock \emph{Mathematical Programming} 80 (1998), 253--264.
\newblock \url{https://doi.org/10.1007/BF01581168}

\bibitem{burger}
A.~P. Burger, E.~J. Cockayne, W.~R. Gr\"undlingh, C.~M. Mynhardt,
J.~H. van Vuuren, and W.~Winterbach,
\newblock Infinite order domination in graphs,
\newblock \emph{Journal of Combinatorial Mathematics and Combinatorial
Computing} 50 (2004), 179--194.
\newblock \url{https://combinatorialpress.com/jcmcc-articles/volume-050/infinite-order-domination-in-graphs/}

\bibitem{ghh}
W.~Goddard, S.~M. Hedetniemi, and S.~T. Hedetniemi,
\newblock Eternal security in graphs,
\newblock \emph{Journal of Combinatorial Mathematics and Combinatorial
Computing} 52 (2005), 169--180.
\newblock \url{https://combinatorialpress.com/jcmcc-articles/volume-052/eternal-security-in-graphs/}

\bibitem{lovasz}
L.~Lov\'asz,
\newblock On the Shannon capacity of a graph,
\newblock \emph{IEEE Transactions on Information Theory} 25 (1979), 1--7.
\newblock \url{https://doi.org/10.1109/TIT.1979.1055985}

\bibitem{mmv}
G.~MacGillivray, C.~M. Mynhardt, and V.~Virgile,
\newblock Eternal domination and clique covering,
\newblock \emph{Electronic Journal of Graph Theory and Applications}
10 (2022), 603--624.
\newblock \url{https://doi.org/10.5614/ejgta.2022.10.2.19}

\bibitem{west}
D.~B. West,
\newblock Eternal domination number,
\newblock Research problems webpage, accessed 26 July 2026.
\newblock \url{https://dwest.web.illinois.edu/regs/eterndom.html}

\end{thebibliography}

\end{document}
